\section{Introduction}\label{sec:introduction}
Network models often combine routing decisions with a choice among services
or operating modes. A route-specific adjustment to a service cost gives a
product of an arc flow and a service weight, and often only a few
route--service pairs carry such an adjustment. Convexifying each product
separately discards restrictions imposed jointly by flow conservation and
the common simplex. The exact joint hull is available through classical
disaggregation over the simplex vertices, but that formulation copies every
arc for every vertex. We ask how the network and the retained products can
reduce this expansion, and when the hull has simple inequalities in the
original coordinates.

We consider a bounded equality-flow polytope
$P=\{x:Ax=b,\ 0\le x\le u\}$ on a directed multigraph and a vector
$y\ge0$ with $\sum_jy_j\le1$. Only products $z_{ej}=x_e y_j$ in a
specified set $\Obs$ are retained. We call these retained product
coordinates \emph{observations} and their joint convex hull $H$. A simplex
label is an index $j$; the unused weight $1-\sum_jy_j$ forms an additional
residual state. The original coordinates are the flows $x$, the weights $y$,
and the retained products $z$. A coefficient is \emph{unit} if it lies in
$\{0,\pm1\}$. Unit statements concern flow and product coefficients (and
auxiliary variables where present); simplex coefficients may depend on the
rational network data.

\paragraph{Known results.}
Disjunctive programming \citep{Balas1998} over the simplex vertices gives
an exact extended formulation of polynomial size, so optimization and
separation over $H$ are already polynomial by linear programming. \citet{DavarniaRichardTawarmalani2017} develop
simultaneous convexification of bilinear functions over a polytope times a
simplex; selecting individual products is a special case.
\citet{KhademniaDavarnia2025}, the closest predecessor, convexify bilinear
terms over network polytopes: their aggregation-and-projection framework
gives a complete original-coordinate description, with explicit tree
inequalities for one simplex variable and a class of forest inequalities for
several. \citet{DavarniaRahimian2026} treat binary simplex variables and
project them out. Our constructions also use shared-simplex gluing
\citep{Davarnia2016}, multiplied-equation elimination
\citep{LibertiPantelides2006}, transportation projections
\citep{KisHorvath2022}, Minkowski-sum support arguments
\citep{GritzmannSturmfels1993,OnnRothblum2004}, and transportation
universality \citep{DeLoeraOnn2006}; \cref{subsec:prior-work} gives the
precise relations. Our aim is to determine what the equality-flow structure
and the pattern of retained products imply \emph{within} the known hull: how
many state coordinates the constructive formulation retains, when a complete separator can use a finite structural library, and
when non-unit product coefficients are unavoidable. Total unimodularity of a
single network does not settle the last question, because sharing aggregate
flows among states and retaining only selected products changes the
projected geometry.

\paragraph{Key idea.}
The relevant quantity is unobserved circulation freedom, not network size.
Circulations separate across the cyclic blocks of the underlying undirected
graph, and labels absent from a block can be merged there. For an observed
label, the remaining freedom is a circulation supported on unobserved arcs.
On a cycle, for example, one observed arc determines the whole state
circulation. A different reduction applies to a chain of parallel arc pairs
with one bypass (a \emph{flat chain}): all serial pairs share one vector of state flows whose
dimension depends on the number of observed labels, not on chain length.

\paragraph{Main results.}
\begin{enumerate}[leftmargin=*,itemsep=5pt]
\item \emph{An observation-sensitive exact formulation.}
For a cyclic block $B$ and a label $j$ observed there, let $\rho_{Bj}$
be the cycle rank of its unobserved subgraph.
\Cref{thm:observed-compression} gives an exact extended formulation with
$\sum_{B,j}\rho_{Bj}$ additional coordinates and unit flow, product, and
auxiliary coefficients in a specified rational row scaling. The count
describes this construction; it is not a minimum over all extended
formulations. It is, however, the minimum number of additional
\emph{individual products} that determine the full state circulation of
every locally observed block--label pair from the observations and the
balance equations on the ambient circulation space
(\cref{prop:minimum-completion}). When all unobserved subgraphs are
forests, the formulation becomes a unit-coefficient description in the
original coordinates on an arbitrary graph (\cref{cor:forest-hull}). The
elimination is the graph case of multiplied-equation elimination
\citep{LibertiPantelides2006}; the contribution is the graph-based count and the
coefficient control of the complete formulation.

\item \emph{Explicit hulls and finite separation procedures.}
Cycle and theta blocks have unit-coefficient interval and support
inequalities, with separation and compact recovery in a linear number of
arithmetic operations (\cref{thm:theta-hull,prop:theta-recovery}). Blocks
consisting of any number of parallel paths have a complete family of subset
inequalities, separated by one maximum-flow computation per block
(\cref{thm:parallel-hull}); this specializes the transportation projection of
\citet{KisHorvath2022} to state slices. For blocks of cycle rank at most $r$,
finite libraries give exact separation in $2^{O(r^2)}N$ rational arithmetic
operations and compact recovery in $2^{O(r^3)}N$, where
$N=|V|+|E|+m+|\Obs|$, after observation grouping
(\cref{thm:bounded-rank,thm:rank-recovery}). On flat chains, whose cycle rank
is unbounded, a fixed number $a$ of observed labels also gives a finite
library, with $O((a+1)L+|\Obs|+m)+2^{O(a^2)}$ operations on a chain of $L$
pairs (\cref{thm:fixed-chain,cor:chain-observed-labels}). None of these
procedures calls an online LP.

\item \emph{Coefficient bounds and their limits.}
At block rank at most $r$, the hull has a description whose flow and
product coefficients are integers of magnitude at most $H_r$ from
\eqref{eq:rank-coefficient-bound}, whose first values are $1,1,2,5$. A unit-capacity $K_4$ with two labels and five products
shows that $H_3=2$ cannot be replaced by one (\cref{cor:rank-three-sharp}).
On a flat chain, every observation pattern with at most two observed labels
has a five-test unit-coefficient oracle, at most three observed labels still
admit unit flow and product coefficients with at most sixteen circuit tests,
and four labels can force a ratio of two
(\cref{thm:two-state-chain,thm:three-state-chain,cor:chain-observed-labels}).
When the number of labels grows, the Fibonacci ratio $F_q$ ($q\ge3$) is
forced with $2q-1$ labels and $5q-4$
observed products, on a flat chain of treewidth two and on simple
series--parallel graphs of maximum degree three
(\cref{thm:fibonacci,cor:fibonacci-simple}). On general networks, every
nonempty bounded rational polytope in a nonnegative orthant is, by a
polynomial-time construction, a coordinate section of a sparse hull with two
labels on an acyclic four-layer network (\cref{thm:universality}). This
transfers the transportation universality of \citet{DeLoeraOnn2006} to
retained product coordinates and forces every integer ratio $M\ge2$ with
unit network data (\cref{cor:universal-coefficients}). These lower bounds
concern actual retained products and persist when affine-hull equations are
added. They do not follow merely from a non-unit aggregation multiplier
\citep[Example~2]{KhademniaDavarnia2025}, from projecting a facet of a larger
product space, or from large coefficients of general $0/1$ polytopes
\citep{AlonVu1997}.
\end{enumerate}

To the best of our knowledge, these formulations, libraries, thresholds,
and constructions are not in the prior work of \cref{subsec:prior-work};
\cref{tab:prior-comparison} separates the established ingredients from the
results claimed here, and \cref{tab:structural-scope} lists the graph
classes.

\begin{table}[tb]\centering\small
\caption{Structural scope of the positive results (unit coefficients as
defined above, in the stated row scaling, before any clearing of simplex
denominators).}\label{tab:structural-scope}
\begin{tabular}{>{\raggedright\arraybackslash}p{.25\linewidth}
                >{\raggedright\arraybackslash}p{.44\linewidth}
                >{\raggedright\arraybackslash}p{.20\linewidth}}\toprule
Structure & Exact representation or procedure & Main locator\\\midrule
Arbitrary equality-flow graph & Observation-sensitive EF with
$\sum_{B,j}\rho_{Bj}$ additional coordinates & \Cref{thm:observed-compression}\\[3pt]
Unobserved arcs form forests in each observed block/state & Unit
original-coordinate formulation & \Cref{cor:forest-hull}\\[3pt]
Cycle or theta blocks & Unit interval/support cuts; compact recovery &
\Cref{thm:theta-hull,prop:theta-recovery}\\[3pt]
Parallel-path blocks & Complete subset cuts; transportation-flow separator &
\Cref{thm:parallel-hull}\\[3pt]
Block cycle rank at most $r$ & Coefficients bounded by $H_r$;
finite support and recovery libraries & \Cref{thm:bounded-rank,thm:rank-recovery}\\[3pt]
Flat chain, $a$ observed labels & Fixed-$a$ profile oracle; unit coefficients
for $a\le3$, sharp at four & \Cref{cor:chain-observed-labels}\\\bottomrule
\end{tabular}
\end{table}

The circulation arguments require equality conservation; if balance
inequalities are converted with slack arcs, all graph assumptions apply to
the transformed network. The general network and block results permit
arbitrary orientations, parallel arcs, loops, zero simplex weights, and
degenerate capacities. The flat-chain
results use the specified directed unit-data topology; they do not assert the
same fixed-label threshold on general nested series--parallel networks.

\paragraph{Computation.}
\Cref{sec:computation} compares exact and numerical implementations on
synthetic instances with full disaggregation, global merging of unused
labels, and elementary one-flow and independent-network baselines. Global
merging is often fastest, and a simple LP beats the exact oracle on some
long-chain queries. When every label is observed somewhere, local compression
reduces 7,215 state-flow variables to 495 (367 after observed elimination)
and the median total time from about 30 to 10 milliseconds. The benefits are
model size and certificates, not a general runtime advantage. A Lean
formalization verifies \cref{prop:disaggregation}, the flat-chain profile
equivalences, and the mathematical claims of \cref{sec:fixed-chains}; the
rest is not formally verified.

\paragraph{Organization.}
\Cref{sec:foundations,sec:blocks} establish the model and the known
disaggregation and gluing facts. \Cref{sec:compression,sec:structured-oracles,sec:bounded-rank}
develop observation-sensitive formulations and graph-structured oracles;
\cref{sec:bounded-rank} also proves the rank-three $K_4$ sharpness result.
\Cref{sec:universality,sec:sp-growth} prove the unbounded coefficient
obstructions: two-label universality and the Fibonacci constructions.
\Cref{sec:fixed-chains} gives the fixed-label chain results, including the
four-label threshold, which reuses an example from \cref{sec:sp-growth}.
\Cref{sec:computation} reports implementations, verification, experiments,
and a two-supplier service-transportation interpretation.
